\subsection{KT-F149-REVIEW: the marked and unmarked statements}
\label{sec:kt-f149-review}

The complete proofs of K17 F149-GENERIC,
F149-MARKED-RECOGNITION, F149-BARE-EXTENSION and
F149-UNMARKED-CURVE have been read in this run. Their locators
are K17.tex, the subsections \emph{Generic recognition and the
necessary profinite input} and \emph{Why the generic marking
cannot simply be omitted}; the principal statements begin at
lines 13654, 13742 and 13908 of the supplied checkpoint-08
source. No fault was found in these arguments under their
stated hypotheses. This is a targeted proof review of those
statements, not an independent proof of the full anabelian
input theorems.

\paragraph{The positive statement and its exact scope.}
In the smooth characteristic-zero situation of F149-RELATIVE,
with normal integral base $S$, generic field $K$ finitely
generated over $\mathbf Q$, and normal line $V$, the checked
criterion is
\[
 s\text{ is induced by }V^\times(S)
 \quad\Longleftrightarrow\quad
 s_K\text{ is induced by }V_K^\times(K).
                                                               \tag{FR.1}
\]
The extension from the generic point is the key algebraic
step: the finite Kummer classes of a rational vector are
unramified at every codimension-one point, so every valuation
of its local coordinate is divisible by all positive integers
and hence zero. Normality extends the vector and its inverse.
The exact commutative-algebra input is EGA~IV$_2$,
Th\'eor\`eme~5.10.5 and its proof, pp.~115--116, as cited
in the supplied proof. Injectivity of inflation along
$G_K\twoheadrightarrow\pi_1(S)$ and injectivity of the full
Kummer map then identify the section and give uniqueness.
The argument does not use finite generation of $\operatorname{Pic}(S)$.

For the marked diagram
\[
 \pi_1(U)\longrightarrow G\longleftarrow G_K
                  \longrightarrow G_k
\]
with identified number field $k$ and regular finitely
generated $K/k$, Mochizuki's field and curve reconstruction
theorems recover $K$ and the generic marked curve.
The source statements are \emph{The local pro-$p$ anabelian
geometry of curves}, Theorems~A and~B (author manuscript
pp.~3--4), Theorem~16.5 together with its following full
profinite comparison (p.~86), and Theorem~17.1 (p.~87).
The statements, the full profinite comparison, and the proof
of Theorem~17.1 were checked directly in the supplied source.
Their regularity and base-Galois hypotheses are present here.
The full proof of Theorem~16.5 was not re-proved in this run.

Once the generic curve has been recovered, a nonzero normal
vector $b$ gives the membership test
\[
 [s_K]-[s_b]\in\kappa_K(K^\times)
       \subset H^1(G_K,\widehat{\mathbf Z}(1)).             \tag{FR.2}
\]
Changing $b$ translates by an element of the same Kummer
image, so the test is independent of that choice. F149-GENERIC
then makes it global. This is a functorial criterion on
realizable full marked diagrams. It is neither a finite
algorithm nor a recognition theorem for arbitrary profinite
groups, and it is not a section-surjectivity theorem.

\paragraph{What the tangential literature already supplies.}
Deligne, \emph{Le groupe fondamental de la droite projective
moins trois points}, \S15.19, p.~239, constructs relative
tangential specialization; \S15.26, p.~242, relates it to
formal coordinates. Stix, \emph{On cuspidal sections of
algebraic fundamental groups}, author version dated
24 February 2012, Proposition~16(2), p.~8, identifies the
cusp-inertia normalizers. His Theorem~17, p.~8, distinguishes
cuspidal packets, and Corollary~32, p.~12, recognizes a
cuspidal section by a cyclotomically normalized inertia
subgroup under the non-Tate hypothesis. These statements
recognize the cusp or its completed packet; they do not
assert that every completed Kummer class comes from an
actual nonzero tangent vector.

The origin of that group-theoretic cusp recognition is
Nakamura's anabelian weight filtration: \emph{Galois rigidity
of the \'etale fundamental groups of punctured projective
lines}, J. Reine Angew. Math.~411 (1990), \S3, and
\emph{Galois rigidity of pure sphere braid groups and
profinite calculus}, J. Math. Sci. Univ. Tokyo~1 (1994),
\S2.1, Theorem~2.1.1 and its proof, pp.~94--97. The latter
statement and complete displayed proof, including
Lemmas~2.1.2 and~2.1.4, were read in the original source;
the former locator is also explicitly given by Stix in \S5. The full
1990 proof was not retrieved and is not claimed checked.

Hoshi--Matsumoto--Nakayama, \emph{Tangential morphisms via
log arithmetic geometry}, RIMS-1995 (June 2026),
Propositions~1--3 and Remark~1, pp.~2--3 of the text
(PDF pp.~3--4), give a log-geometric formulation, coordinate
independence up to homotopy, and composition. In particular,
Remark~1 makes dependence on the normal cotangent basis
explicit. This supports the functoriality used by F149; it
does not remove the generic marking in (FR.2).

Sa\"idi's \emph{Good Sections of Arithmetic Fundamental
Groups}, arXiv:1010.1313v1, Theorems~2.6--2.7,
pp.~33--39, treats lifting uniformly good sections to
cuspidally abelian groups over slim, regular fields.
This is a different lifting problem. The published related
paper is \emph{The cuspidalisation of sections of arithmetic
fundamental groups}, Adv. Math.~230 (2012), 1931--1954;
its abstract was checked, but no identification of its
statement numbers with the preprint is asserted here.

\paragraph{The negative statement survives a full read.}
F149-UNMARKED-CURVE uses
$S=\mathbf G_{m,\mathbf Q}$ and
$U=S\times(\mathbf P^1\setminus\{0,1,\infty\})$.
Scaling the geometric base group
$\widehat{\mathbf Z}(1)$ by a profinite unit $u$, with
$u_2=3$ and $u_\ell=1$ for $\ell\ne2$, extends to the
full relative curve diagram and preserves its arithmetic
quotient and cusp data. It carries the tangent section of
$t b$ to slope $u^{-1}$ on base inertia. Actual global units
are $a t^m$ with $m\in\mathbf Z$, and their slopes are
integers; $u^{-1}$ has two-adic component $1/3$ and is not
an integer. The proof checks conjugacy as well: inertia
coboundaries vanish on geometric base inertia.

Thus a reference tangent vector and a rational base section
do not fix the missing marking when automorphisms may move
$G$. This does not conflict with Mochizuki's full-faithfulness
theorem, because that theorem retains the appropriate base
field Galois marking. No counterexample with $G$ fixed
pointwise, with its scheme realization fixed, or with $G_K$
retained is claimed.

\paragraph{Attribution and novelty.}
The anabelian reconstruction, tangent functors, cusp
normalizers and completed Kummer packets are established
literature. The K17 contribution reviewed here is their
assembly into the exact generic-to-global test (FR.1)--(FR.2)
and the explicit smaller-diagram automorphism obstruction.
No exhaustive priority claim for those formulations has
been established by this literature search. In Grothendieck's
terms, the result addresses the distinction between the
geometric subtorsor and the full profinite torsor on
cote~149 scan~85, with precisely the additional generic
marking of scans~83--84.

